% (C) 2026 Ralf Stephan, in collaboration with Claude Code.
% Released under CC0 1.0 Universal (public-domain dedication).
% See https://creativecommons.org/publicdomain/zero/1.0/

\documentclass[11pt]{article}
\usepackage[T1]{fontenc}
\usepackage[utf8]{inputenc}
\usepackage{lmodern}
\usepackage[margin=2.5cm]{geometry}
\usepackage{amsmath,amssymb}
\usepackage{braket}
\usepackage[hidelinks]{hyperref}

% Lean names and paths: verbatim, breakable at . _ /
\DeclareUrlCommand\lean{\urlstyle{tt}}
\def\UrlBreaks{\do\.\do\_\do\/}

\title{Open problems:\\ Non-uniqueness of the moment problem for Schwartz functions}
\author{}
\date{}

\begin{document}
\maketitle

\noindent
Open problems raised by the technical brief in this directory, and by its Lean formalization
(see \lean{formalization-log.tex}).

\section*{From the brief}
\begin{enumerate}
\item \textbf{Moments see only distances to $0$.} Although the topology of
  $\mathcal{S}(\mathbb{R})$ is generated by the seminorms $\|f\|^{L^2}_{\alpha,\beta}$, ``the
  moments, then, can only give information about all distances to 0, and not to all other
  functions''. This explanation of the failure is not made precise. (The seminorm equivalence
  itself is formalized for the brief ``Equivalence of sup-norm, $L^1$, and $L^2$ seminorms on the
  Schwartz space''.) (Brief, Remark after Theorem~1; log, ``Differences from the brief'': not
  formalized.)
\item \textbf{Locality as the obstruction.} ``The key obstruction is that polynomials of $X$ and
  $P$ are local operators and can never detect phase differences between spatially separated
  parts.'' The footnote ties locality to Peetre's theorem. The general statement, that
  $\braket{f|A|g} = 0$ for every local operator $A$ and disjointly supported $f, g$, so that no
  such $A$ separates $\psi_\theta$ from $\psi_{\theta'}$, is not formalized. (Brief, Sec.~3 and
  footnote~1; log, ``Differences from the brief'': not formalized.)
\item \textbf{Other state spaces.} ``Consequently any space that contains functions with compact
  support will suffer the same consequences.'' Not stated or proved beyond
  $\mathcal{S}(\mathbb{R})$, e.g.\ for $\mathbb{R}^n$ or other domains. (Brief, Sec.~3.)
\item \textbf{Which states are determined by their moments.} The brief asks, ``As in classical
  probability'', whether each state can be identified by its moments of polynomials in $X$ and
  $P$, and answers only negatively. Conditions under which the moments do determine the state,
  the analogue of classical determinacy criteria, are not addressed. (Brief, Sec.~1.)
\end{enumerate}

\section*{From the formalization}
\begin{enumerate}
\item None found.
\end{enumerate}

\end{document}
